\documentclass[11pt,a4paper]{article}

\usepackage[T1]{fontenc}
\usepackage[utf8]{inputenc}
\usepackage{lmodern}
\usepackage[a4paper,top=2.4cm,bottom=2.2cm,left=2.2cm,right=2.2cm]{geometry}
\usepackage{amsmath}
\usepackage{amssymb}
\usepackage{graphicx}
\graphicspath{{./}}
\usepackage{array}
\usepackage{makecell}
\renewcommand{\arraystretch}{1.25}
\usepackage{enumitem}
\usepackage{fancyhdr}
\usepackage{xcolor}

\newcommand{\ohms}{\ensuremath{\Omega}}
\newcommand{\degC}{\ensuremath{^{\circ}}C}
\newcommand{\dg}{\ensuremath{^{\circ}}}
\newcommand{\pow}[1]{\ensuremath{^{#1}}}

\newcommand{\topiclabel}[1]{{\hfill\normalfont\small\itshape\textcolor{black!55}{#1}}}

% ===== Per-question head: \examq{paper-id}{number}{topics}{marks} =====
% Same call in the question paper and the mark scheme, so the two join on
% paper-id + number. topics = middle-dot-separated list (\textperiodcentered
% between topics; the dot never appears inside a single topic name). marks is
% stored here so \total prints it with no argument (one source of truth).
\newcommand{\swmarks}{}
\newcommand{\examq}[4]{%
  \gdef\swmarks{#4}%
  \par\addvspace{16pt}\noindent
  {\large\bfseries #1 - Q#2 - #3 - #4}\par\vspace{5pt}%
}

\newlength{\anslineskip}
\setlength{\anslineskip}{8mm}

\newcounter{ansln}
\newcommand{\Alines}[2]{%
  \par\vspace{1mm}%
  \setcounter{ansln}{1}%
  \loop\ifnum\value{ansln}<#1\relax
    \noindent\mbox{}\dotfill\par\vspace{\anslineskip}%
    \stepcounter{ansln}%
  \repeat
  \noindent\mbox{}\dotfill\hspace{0.6em}\mbox{#2}\par
  \vspace{1mm}%
}

\newcommand{\plainline}{\par\noindent\mbox{}\dotfill\par\vspace{\anslineskip}}

\newcommand{\labelline}[1]{\par\vspace{1mm}\noindent #1\hspace{0.5em}\dotfill\par\vspace{\anslineskip}}

\newcommand{\ansval}[3]{%
  \par\vspace{5mm}%
  \noindent\hspace*{3.5cm}#1\hspace{0.5em}\dotfill\hspace{0.5em}#2\hspace{0.5em}\makebox[1.1cm][r]{#3}\par
  \vspace{3mm}%
}

\renewcommand{\markright}[1]{\par\nopagebreak\noindent\hspace*{\fill}\mbox{#1}\par\vspace{1mm}}

\newcommand{\total}{\par\vspace{2mm}\noindent\hspace*{\fill}[Total:\ \swmarks]\par\vspace{2mm}}

\newcommand{\qfig}[3][0.5\linewidth]{%
  \par\vspace{2mm}
  \begin{center}%
    \IfFileExists{#2}%
      {\includegraphics[width=#1]{#2}}%
      {\fbox{\parbox[c][26mm][c]{#1}{\centering\ttfamily\footnotesize #2}}}\\[4pt]
    \textbf{#3}%
  \end{center}%
  \vspace{2mm}\par
}

\newlist{parts}{enumerate}{1}
\setlist[parts]{label=\textbf{(\alph*)},leftmargin=2.1em,labelsep=0.6em,itemsep=10pt,topsep=6pt,parsep=4pt}

\newlist{subparts}{enumerate}{1}
\setlist[subparts]{label=\textbf{(\roman*)},leftmargin=2.3em,labelsep=0.6em,itemsep=10pt,topsep=6pt,parsep=4pt}

\pagestyle{fancy}
\fancyhf{}
\fancyhead[L]{\footnotesize 4024/22/M/J/25}
\fancyhead[C]{\footnotesize Cambridge O Level Mathematics (Syllabus D) -- Paper 2}
\fancyhead[R]{\footnotesize May/June 2025}
\fancyfoot[C]{\footnotesize\thepage}
\renewcommand{\headrulewidth}{0.4pt}
\renewcommand{\footrulewidth}{0pt}

\setlength{\parindent}{0pt}

\begin{document}

\examq{4024/22/M/J/25}{1}{Geometric Fundamentals}{1}
\qfig[0.5\linewidth]{4024_22_M_J_25_fig1.png}{}
The diagram shows part of a pattern with rotational symmetry of order 4.

Complete the diagram.
\markright{[1]}
\total

\examq{4024/22/M/J/25}{2}{Number Fundamentals}{3}
\begin{parts}
  \item Write eighteen thousand and twelve in figures.
  \ansval{}{}{[1]}

  \item Write down a prime number between 20 and 30.
  \ansval{}{}{[1]}

  \item Find the reciprocal of $\dfrac{1}{9}$\,.
  \ansval{}{}{[1]}
\end{parts}
\total

\examq{4024/22/M/J/25}{3}{Ratio \& Proportion}{4}
\begin{parts}
  \item Write the ratio $\;175\ \mathrm{ml} : 2.5\ \mathrm{litres}\;$ in its simplest form.
  \vspace{3cm}
  \par\vspace{5mm}\noindent\hspace*{\fill}\makebox[3cm]{\dotfill} : \makebox[3cm]{\dotfill}\hspace{0.6em}\mbox{[2]}\par\vspace{3mm}

  \item Dan, Erika and Fatik share \$540 in the ratio $\;4:5:3\;$.

  Find the amount that Erika receives.
  \vspace{3.5cm}
  \ansval{\$}{}{[2]}
\end{parts}
\total

\examq{4024/22/M/J/25}{4}{Angles in Polygons \& Parallel Lines}{2}
\qfig[0.75\linewidth]{4024_22_M_J_25_fig2.png}{}
Find the value of $x$ and the value of $y$.
\ansval{$x =$}{}{}
\ansval{$y =$}{}{}
\markright{[2]}
\total

\examq{4024/22/M/J/25}{5}{Geometric Fundamentals}{1}
Convert $6.1\,\mathrm{m}^{2}$ to $\mathrm{cm}^{2}$.
\vspace{2cm}
\ansval{}{cm\pow{2}}{[1]}
\total

\examq{4024/22/M/J/25}{6}{Probability Fundamentals}{3}
In a survey, students are asked to choose their favourite type of movie.\\
The table shows the results.

\begin{center}
\begin{tabular}{|l|c|c|c|c|}
\hline
Type of movie & Romance & Thriller & Science Fiction & Other \\
\hline
Relative frequency & 0.1 & & 0.35 & 0.3 \\
\hline
\end{tabular}
\end{center}

\begin{parts}
  \item Calculate the relative frequency of a student choosing Thriller.
  \vspace{1cm}
  \ansval{}{}{[2]}

  \item 500 students take part in the survey.

  Calculate the number of students who choose Science Fiction.
  \vspace{1cm}
  \ansval{}{}{[1]}
\end{parts}
\total

\examq{4024/22/M/J/25}{7}{Time, Currency \& Conversions}{3}
Chris wants to exchange \$350 for euros (\texteuro) at the bank.\\
The bank only gives euros in multiples of \texteuro5.\\
The exchange rate is \$1 $=$ \texteuro0.92\,.

Calculate the number of euros he receives and his change from \$350.
\vspace{6cm}
\ansval{Chris receives \texteuro}{}{}
\ansval{His change is \$}{}{}
\markright{[3]}
\total

\examq{4024/22/M/J/25}{8}{Angles in Polygons \& Parallel Lines}{2}
Calculate the interior angle of a regular octagon.
\vspace{4cm}
\ansval{}{}{[2]}
\total

\examq{4024/22/M/J/25}{9}{Set Notation \& Venn Diagrams}{3}
\begin{parts}
  \item On the Venn diagram, shade the region represented by $\;(A\cap B)'\;$.
  \qfig[0.55\linewidth]{4024_22_M_J_25_fig3.png}{}
  \markright{[1]}

  \item In a class of 22 students:
  \begin{itemize}[leftmargin=2.2em,itemsep=2pt,topsep=4pt]
    \item 12 play the piano
    \item 9 play the guitar
    \item 6 do \textbf{not} play the piano or the guitar.
  \end{itemize}

  Find the number of students who play both the piano and the guitar.

  You may use the Venn diagram to help you.
  \qfig[0.55\linewidth]{4024_22_M_J_25_fig4.png}{}
  \ansval{}{}{[2]}
\end{parts}
\total

\examq{4024/22/M/J/25}{10}{Algebraic Roots \& Indices}{2}
Simplify.
\[\dfrac{7x^{2}y^{5}}{x^{3}y^{2}}\]
\vspace{3cm}
\ansval{}{}{[2]}
\total

\examq{4024/22/M/J/25}{11}{Simple \& Compound Interest, Growth \& Decay}{6}
\begin{parts}
  \item Zaya invests \$4500 in a savings account.\\
  The account pays \textbf{simple interest} at a rate of 3.2\% per year.

  Calculate the value of the investment at the end of 5 years.
  \vspace{4.5cm}
  \ansval{\$}{}{[3]}

  \item Aisha invests \$2750 in a different savings account.\\
  This account pays \textbf{compound interest} at a rate of 2.1\% per year.

  Calculate the total amount of interest she receives at the end of 3 years.
  \vspace{4.5cm}
  \ansval{\$}{}{[3]}
\end{parts}
\total

\examq{4024/22/M/J/25}{12}{Expanding \& Factorising Brackets}{2}
Factorise.
\[7hx+6fy-21fx-2hy\]
\vspace{4cm}
\ansval{}{}{[2]}
\total

\examq{4024/22/M/J/25}{13}{Probability Diagrams}{3}
A bag contains 9 red counters and 3 green counters.\\
Lee takes two counters from the bag at random without replacement.
\begin{parts}
  \item Complete the tree diagram.
  \qfig[0.85\linewidth]{4024_22_M_J_25_fig5.png}{}
  \markright{[2]}

  \item Work out the probability that both counters are red.
  \vspace{1.5cm}
  \ansval{}{}{[1]}
\end{parts}
\total

\examq{4024/22/M/J/25}{14}{Circle Theorems}{3}
\qfig[0.72\linewidth]{4024_22_M_J_25_fig6.png}{}
$PQRS$ is a cyclic quadrilateral.\\
$APB$ is a tangent to the circle at $P$.\\
Angle $SPB = 37\dg$ and angle $PSQ = 85\dg$.
\begin{parts}
  \item Find angle $PQS$.
  \ansval{Angle $PQS =$}{}{[1]}

  \item Find angle $QRS$.
  \vspace{1.5cm}
  \ansval{Angle $QRS =$}{}{[2]}
\end{parts}
\total

\examq{4024/22/M/J/25}{15}{Sequences}{7}
These are the first four patterns in a sequence made from black beads and from white beads.
\qfig[0.85\linewidth]{4024_22_M_J_25_fig7.png}{}
\begin{parts}
  \item Complete the table for Pattern 5.

  \begin{center}
  \begin{tabular}{|l|c|c|c|c|c|}
  \hline
  Pattern number & 1 & 2 & 3 & 4 & 5 \\
  \hline
  Number of white beads & 5 & 7 & 9 & 11 & \\
  \hline
  Total number of beads & 6 & 10 & 15 & 21 & \\
  \hline
  \end{tabular}
  \end{center}
  \markright{[1]}

  \item Find an expression for
  \begin{subparts}
    \item the number of white beads in Pattern $n$
    \vspace{2cm}
    \ansval{}{}{[2]}

    \item the total number of beads in Pattern $n$.
    \vspace{4cm}
    \ansval{}{}{[2]}
  \end{subparts}

  \item Bill has only 88 white beads but lots of black beads.\\
  Pattern $p$ is the largest possible pattern he can make using these beads.

  Work out the value of $p$.
  \vspace{3.5cm}
  \ansval{$p =$}{}{[2]}
\end{parts}
\total

\examq{4024/22/M/J/25}{16}{Powers, Roots \& Standard Form}{2}
The population of Kenya is $5.71\times10^{7}$ people.\\
The area of Kenya is $6\times10^{5}\,\mathrm{km}^{2}$.\\
The population density is the number of people per $\mathrm{km}^{2}$.

Calculate the population density of Kenya.
\vspace{4cm}
\ansval{}{people\,/\,km\pow{2}}{[2]}
\total

\examq{4024/22/M/J/25}{17}{Functions}{3}
\begin{parts}
  \item \[\mathrm{f}(x) = \dfrac{10-x}{3}\]
  Find $\;\mathrm{f}(-8)\;$.
  \vspace{2.5cm}
  \ansval{}{}{[1]}

  \item \[\mathrm{g}(x) = 4x+3\]
  Find $\;\mathrm{g}^{-1}(x)\;$.
  \vspace{3.5cm}
  \ansval{$\mathrm{g}^{-1}(x) =$}{}{[2]}
\end{parts}
\total

\examq{4024/22/M/J/25}{18}{Averages \& Range \textperiodcentered\ Histograms}{7}
300 students take part in a competition.\\
The table shows information about their scores.

\begin{center}
\small
\begin{tabular}{|l|c|c|c|c|c|}
\hline
Score ($s$) & $0 < s \leqslant 30$ & $30 < s \leqslant 40$ & $40 < s \leqslant 60$ & $60 < s \leqslant 70$ & $70 < s \leqslant 100$ \\
\hline
Frequency & 24 & 70 & 88 & 76 & 42 \\
\hline
\end{tabular}
\end{center}

\begin{parts}
  \item Calculate an estimate of the mean score.
  \vspace{5cm}
  \ansval{}{}{[4]}

  \item \leavevmode
  \qfig[0.8\linewidth]{4024_22_M_J_25_fig8.png}{}
  On the grid, complete the histogram to represent the data in the table.
  \markright{[3]}
\end{parts}
\total

\examq{4024/22/M/J/25}{19}{Curve Sketching \& Gradient Analysis}{8}
\begin{parts}
  \item Complete the table of values for $\;y = x^{3}-5x+4$.

  \begin{center}
  \begin{tabular}{|c|c|c|c|c|c|c|c|}
  \hline
  $x$ & $-3$ & $-2$ & $-1$ & 0 & 1 & 2 & 3 \\
  \hline
  $y$ & & 6 & 8 & 4 & 0 & 2 & 16 \\
  \hline
  \end{tabular}
  \end{center}
  \markright{[1]}

  \item Draw the graph of $\;y = x^{3}-5x+4\;$ for $\;-3 \leqslant x \leqslant 3$.
  \qfig[0.85\linewidth]{4024_22_M_J_25_fig9.png}{}
  \markright{[4]}

  \item By drawing a suitable straight line on the grid, find the solutions of the equation $\;x^{3}-5x-3 = 0$.
  \par\vspace{6mm}\noindent\hspace*{\fill}$x =$ \makebox[2.6cm]{\dotfill} , $x =$ \makebox[2.6cm]{\dotfill} , $x =$ \makebox[2.6cm]{\dotfill}\hspace{0.6em}\mbox{[3]}\par\vspace{3mm}
\end{parts}
\total

\examq{4024/22/M/J/25}{20}{Quadratic Equations \textperiodcentered\ Forming \& Solving Equations \textperiodcentered\ Area \& Perimeter}{9}
$ABCD$ and $EFGD$ are rectangles.
\qfig[0.8\linewidth]{4024_22_M_J_25_fig10.png}{}
\begin{parts}
  \item \leavevmode
  \begin{subparts}
    \item Show that $\;CD = x-1\;$.
    \vspace{2.5cm}
    \markright{[1]}

    \item The area of rectangle $ABCD$ is $\dfrac{1}{5}$ of the area of rectangle $EFGD$.

    Form an equation in $x$ and show that it simplifies to $\;x^{2}-6x+3 = 0\;$.
    \vspace{7cm}
    \markright{[3]}
  \end{subparts}

  \item \leavevmode
  \begin{subparts}
    \item Use the quadratic formula to solve the equation $\;x^{2}-6x+3 = 0$.\\
    You must show all your working and give your answers correct to 2 decimal places.
    \vspace{7cm}
    \par\vspace{6mm}\noindent\hspace*{\fill}$x =$ \makebox[3.2cm]{\dotfill} or $x =$ \makebox[3.2cm]{\dotfill}\hspace{0.6em}\mbox{[3]}\par\vspace{3mm}

    \item Calculate the shaded area.
    \vspace{5cm}
    \ansval{}{units\pow{2}}{[2]}
  \end{subparts}
\end{parts}
\total

\examq{4024/22/M/J/25}{21}{Volume \& Surface Area}{11}
\qfig[0.6\linewidth]{4024_22_M_J_25_fig11.png}{}
The diagram shows a solid formed by joining a cone to a hemisphere.\\
The diameter of the cone is 14\,cm and the diameter of the hemisphere is 14\,cm.\\
The total height of the solid is 24\,cm.
\begin{parts}
  \item Calculate the volume of the solid.
  \vspace{6cm}
  \ansval{}{cm\pow{3}}{[3]}

  \item Show that the total surface area of the solid is $712\,\mathrm{cm}^{2}$, correct to the nearest integer.
  \vspace{9cm}
  \markright{[5]}

  \item A smaller solid is mathematically similar to this solid.\\
  The total surface area of the smaller solid is $242\,\mathrm{cm}^{2}$.

  Calculate the total height of the smaller solid.
  \vspace{5cm}
  \ansval{}{cm}{[3]}
\end{parts}
\total

\examq{4024/22/M/J/25}{22}{Sine, Cosine Rule \& Area of Triangles}{9}
\qfig[0.7\linewidth]{4024_22_M_J_25_fig12.png}{}
The diagram shows a triangular field $ABC$.\\
$AB = 150$\,m, $\;BC = 187$\,m and angle $ABC = 112\dg$.
\begin{parts}
  \item Fencing is needed for the perimeter of the field.\\
  Fencing is sold in rolls of length 20\,m.

  Calculate the number of rolls of fencing needed.
  \vspace{7cm}
  \ansval{}{}{[5]}

  \item Calculate the shortest distance from $B$ to $AC$.
  \vspace{6cm}
  \ansval{}{m}{[4]}
\end{parts}
\total

\examq{4024/22/M/J/25}{23}{Compound Measures \textperiodcentered\ Rounding, Estimation \& Bounds}{3}
Zara cycles a distance of 500 metres, correct to the nearest 5 metres.\\
This takes 24.7 seconds, correct to the nearest 0.1 seconds.

Calculate the lower bound of her average speed.
\vspace{7cm}
\ansval{}{m/s}{[3]}
\total

\examq{4024/22/M/J/25}{24}{Algebraic Fractions}{3}
Express as a single fraction in its simplest form.
\[\dfrac{5}{2x+1}-\dfrac{2}{4x-3}\]
\vspace{7cm}
\ansval{}{}{[3]}
\total

\end{document}